% Appendix E source fragment. Do not input this file into the six-page paper.
% The overall portfolio wrapper should provide the Appendix E title page and
% appendix numbering before inputting this content.
\section{Testing and Supplementary Results}
\label{app:testing-results}

\subsection{Evaluation Definitions}
For a dB response with $N$ evaluated values, the reported error metrics are
\begin{align}
    \operatorname{MAE}_{\mathrm{dB}}
      &=\frac{1}{N}\sum_{i=1}^{N}|y_i-\widehat{y}_i|, \\
    \operatorname{RMSE}_{\mathrm{dB}}
      &=\sqrt{\frac{1}{N}\sum_{i=1}^{N}(y_i-\widehat{y}_i)^2}.
    \label{eq:appendix-db-errors}
\end{align}
The deep-null subset is
$\mathcal{D}=\{i:y_i\geq69.189752\ \mathrm{dB}\}$, where the threshold is the
99th percentile of the pooled six-path training targets. The common-target test
set contains 3,310 such $IL_{7,1}$ design--frequency points, and
\begin{equation}
    \operatorname{DeepNullMAE}
      =\frac{1}{|\mathcal{D}|}
       \sum_{i\in\mathcal{D}}|y_i-\widehat{y}_i|.
    \label{eq:appendix-deep-null}
\end{equation}
This S7-only quantity must not be confused with the native six-path deep-null
values stored by the Curve and Full-S-Matrix runs.

For complete complex matrices, with each $S$-parameter counted once per design,
frequency and port pair,
\begin{align}
    \operatorname{ComplexMAE}
      &=\frac{1}{N}\sum_{i=1}^{N}
        |\widehat{S}_i-S_i|, \\
    \operatorname{ComplexNRMSE}
      &=\sqrt{\frac{\sum_i|\widehat{S}_i-S_i|^2}
                         {\sum_i|S_i|^2}}.
    \label{eq:appendix-complex-errors}
\end{align}
These dimensionless complex-domain metrics are not numerically comparable with
dB error.

The reciprocity residual is the normalised Frobenius norm
$\|\widehat{\mathbf S}-\widehat{\mathbf S}^{\mathsf T}\|_F/
\|\widehat{\mathbf S}\|_F$. Sampled passivity requires
$\sigma_{\max}(\widehat{\mathbf S})\leq1$ at each design and frequency. The
causality result is a normalised Hilbert-transform residual over only the
available 0.5--100~GHz positive-frequency band; it is a relative diagnostic,
not a full-band causality test.

\subsection{Selected-Model Validation and Test Errors}
Table~\ref{tab:appendix-s7-errors} gives the complete common-target MAE and RMSE
record. All values were recomputed by NB07 from the runs named in
Table~\ref{tab:appendix-selected-runs}.

\begin{table*}[t]
    \caption{Selected-model $IL_{7,1}$ validation and test errors (dB).}
    \label{tab:appendix-s7-errors}
    \centering
    \begin{tabular}{@{}lrrrr@{}}
        \toprule
        & \multicolumn{2}{c}{Validation} & \multicolumn{2}{c}{Test} \\
        Model & MAE & RMSE & MAE & RMSE \\
        \midrule
        Scalar Ridge     &  7.3237 & 10.8112 &  7.5051 & 12.3342 \\
        Vector Ridge     &  7.3237 & 10.8112 &  7.5051 & 12.3342 \\
        Polynomial Ridge &  7.2987 & 10.7929 &  7.4968 & 12.3411 \\
        Random Forest    &  7.6717 & 11.4302 &  7.8043 & 12.6891 \\
        Neural MLP       &  7.3296 & 10.8119 &  7.5073 & 12.3324 \\
        Polynomial MLP   &  7.3183 & 10.8090 &  7.5087 & 12.3425 \\
        Curve Neural     &  7.1679 & 10.7409 &  7.3883 & 12.3097 \\
        Full S-Matrix    & 10.6129 & 13.0990 & 10.7903 & 14.3446 \\
        \bottomrule
    \end{tabular}
\end{table*}

\begin{table*}[t]
    \caption{Common-target $IL_{7,1}$ test MAE and DeepNullMAE (dB).}
    \label{tab:appendix-deep-null-errors}
    \centering
    \begin{tabular}{@{}lrr@{}}
        \toprule
        Predictor & MAE & DeepNullMAE \\
        \midrule
        Global training mean & 12.3535 & 72.4318 \\
        Training mean curve  &  7.4254 & 62.0567 \\
        Scalar Ridge         &  7.5051 & 61.6731 \\
        Vector Ridge         &  7.5051 & 61.6731 \\
        Polynomial Ridge     &  7.4968 & 62.1812 \\
        Random Forest        &  7.8043 & 60.6063 \\
        Neural MLP           &  7.5073 & 61.7068 \\
        Polynomial MLP       &  7.5087 & 62.1162 \\
        Curve Neural         &  7.3883 & 62.3265 \\
        Full S-Matrix        & 10.7903 & 63.0649 \\
        \bottomrule
    \end{tabular}
\end{table*}

\subsection{Paired Design-Level Comparisons}
For each transition, NB07 first averages absolute $IL_{7,1}$ error over the 200
frequencies of each test design. It then forms the paired difference
$d_n=\operatorname{MAE}_{n,\mathrm{current}}-
\operatorname{MAE}_{n,\mathrm{previous}}$. Five thousand percentile-bootstrap
samples resample the 1,406 complete test designs with replacement using seed
7135. Negative differences favour the current model. The intervals are
conditional on the fixed saved models and do not include retraining,
hyperparameter-search or repeated-test-exposure uncertainty.

\begin{table*}[t]
    \caption{Paired test-design bootstrap comparisons.}
    \label{tab:appendix-bootstrap}
    \centering
    \scriptsize
    \begin{tabular}{@{}lrrrl@{}}
        \toprule
        Transition & $\Delta$ MAE (dB) & Change & 95\% interval (dB)
                   & Interpretation \\
        \midrule
        Mean curve $\rightarrow$ Scalar Ridge
          & +0.0797 & +1.07\% & [+0.0515,+0.1081] & directional \\
        Scalar $\rightarrow$ Vector Ridge
          & +0.0000 & +0.00\% & [+0.0000,+0.0000] & equivalent \\
        Vector $\rightarrow$ Polynomial Ridge
          & -0.0082 & -0.11\% & [-0.0348,+0.0195] & equivalent \\
        Polynomial Ridge $\rightarrow$ Random Forest
          & +0.3074 & +4.10\% & [+0.2020,+0.4082] & meaningful increase \\
        Random Forest $\rightarrow$ Neural MLP
          & -0.2969 & -3.80\% & [-0.4013,-0.1884] & meaningful reduction \\
        Vector Ridge $\rightarrow$ Neural MLP
          & +0.0022 & +0.03\% & [-0.0051,+0.0095] & equivalent \\
        Neural MLP $\rightarrow$ Polynomial MLP
          & +0.0014 & +0.02\% & [-0.0307,+0.0339] & equivalent \\
        Polynomial MLP $\rightarrow$ Curve Neural
          & -0.1204 & -1.60\% & [-0.1613,-0.0774] & directional \\
        Curve Neural $\rightarrow$ Full S-Matrix
          & +3.4020 & +46.05\% & [+3.3104,+3.4901] & meaningful increase \\
        \bottomrule
    \end{tabular}
\end{table*}

The practical-equivalence band was predeclared as $[-0.10,+0.10]$~dB.
``Directional'' means the interval excludes zero but crosses one boundary of
that band. ``Meaningful increase/reduction'' means the whole interval lies
outside the band.

\subsection{Complete-Matrix Accuracy and Physical Diagnostics}
\begin{table*}[t]
    \caption{Complete-matrix Complex MAE and Complex NRMSE.}
    \label{tab:appendix-complex-errors}
    \centering
    \begin{tabular}{@{}llrr@{}}
        \toprule
        Predictor & Split & Complex MAE & Complex NRMSE \\
        \midrule
        Training-mean matrix & validation & 0.081178 & 0.837008 \\
                             & test       & 0.080399 & 0.834299 \\
        Full S-Matrix Neural & validation & 0.069033 & 0.790785 \\
                             & test       & 0.068567 & 0.788925 \\
        \bottomrule
    \end{tabular}
\end{table*}

\begin{table*}[t]
    \caption{Complete-matrix physical diagnostics.}
    \label{tab:appendix-physics}
    \centering
    \scriptsize
    \begin{tabular}{@{}llrrrr@{}}
        \toprule
        Split & Matrix & Reciprocity & Passivity violations
              & Maximum $\sigma$ & Causality residual \\
        \midrule
        validation & truth      & $3.41\times10^{-11}$ & 0/281,200
                   & 0.999998 & 0.318476 \\
        validation & prediction & 0 & 0/281,200
                   & 0.924104 & 0.481261 \\
        test & truth      & $7.00\times10^{-11}$ & 0/281,200
             & 0.999978 & 0.317284 \\
        test & prediction & 0 & 0/281,200
             & 0.914132 & 0.481240 \\
        \bottomrule
    \end{tabular}
\end{table*}

Only reciprocity is guaranteed, because it follows from matrix reconstruction.
The passivity counts cover the sampled design--frequency matrices only. The
finite positive-frequency causality residual cannot establish causality over an
unbounded, two-sided spectrum.

\subsection{Selected Runs and Reproduction Record}
\begin{table*}[t]
    \caption{Selected runs loaded by NB07.}
    \label{tab:appendix-selected-runs}
    \centering
    \scriptsize
    \begin{tabular}{@{}ll@{}}
        \toprule
        Model & Run identifier \\
        \midrule
        Scalar Ridge     & \nolinkurl{20260804T143344Z_scalar_ridge} \\
        Vector Ridge     & \nolinkurl{20260804T143708Z_vector_ridge} \\
        Polynomial Ridge & \nolinkurl{20260804T143713Z_polynomial_ridge} \\
        Random Forest    & \nolinkurl{20260804T143727Z_random_forest} \\
        Neural MLP       & \nolinkurl{20260804T144030Z_neural_mlp} \\
        Polynomial MLP   & \nolinkurl{20260804T144859Z_polynomial_neural_mlp} \\
        Curve Neural     & \nolinkurl{20260808T125830Z_curve_neural} \\
        Full S-Matrix    & \nolinkurl{20260804T152940Z_full_smatrix_neural} \\
        \bottomrule
    \end{tabular}
\end{table*}

NB07 loads these entries from \nolinkurl{outputs/models/selected.json}, checks
that every artifact exists, verifies benchmark run identifiers, aligns the
design-level splits, and recomputes saved metrics. Its final audit performed 80
metric-reproduction checks: all passed within their recorded tolerances. No
model was trained, promoted or selected in NB07.

The reproducible evaluation inputs are the named dataset, the project seed 128,
the 4,218/1,406/1,406 design split, the eight selected run manifests and the
paired NB07 source/notebook. Earlier NB03--NB06 development inspected test
results, so this record supports retrospective held-out evaluation rather than
an untouched final confirmation.

The machine-readable evidence exported by NB07 is stored in
\nolinkurl{project_portfolio/evidence/}: the common-target and complete-matrix
tables, all paired transitions, physical diagnostics, selected-run provenance,
80 metric checks, and the plotted fixed-design response data. The paper figure
is regenerated from \nolinkurl{response_design_5491.csv}; it is not a manually
selected response.
